\section*{Cluster A Step 21: Mode-B Unifying Carrier}

\paragraph{Grade.}
This is a finite Mode-B generation result. The carrier is a minimal operational kernel; it is not a physical unification theory and it does not certify frame transfer.

\paragraph{Carrier.}
The kernel uses five finite slots split by the SM-side sector lens \(q\), bridge operations that \(q\) forgets, a primitive trace-neutral integer generator, a trace metric, and dual/pair rewrite operations. The promoted object \(U\) is non-descending: two promoted bridge states have the same \(q\)-readout and different bridge readouts.

\paragraph{Audited descent.}
The audited expression \(C(U)\) descends through \(q\). The primitive trace audit computes
\[
  \mathrm{Tr}(Y^2)=\frac{5}{6},\qquad
  \mathrm{Tr}(T_3^2)=\frac{1}{2},
\]
so the generated normalization is
\[
  \frac{\mathrm{Tr}(Y^2)}{\mathrm{Tr}(T_3^2)}=\frac{5}{3}
\]
and
\[
  \sin^2\theta_W=\frac{1}{1+5/3}=\frac{3}{8}.
\]
The charge unit is \(1/6\). This recognition-lands on the SU(5)-shaped value after the audit; it is not imported as a primitive.

\paragraph{No-smuggling gates.}
The primitive-exclusion, dependency-trace, ablation, negative-control, Stage-II, and no-single-axiom-equivalence gates all pass. Stage II reproduces zero mixed color-charge, mixed weak-charge, cubic-charge, and charge-gravity anomaly sums from the generated dual/pair rewrite content.

\paragraph{Boundary.}
The result is finite-kernel evidence that the framework grammar can generate a GUT-shaped normalization relation without using it as a primitive. It is not a claim that physical unification is closed; the Step-20 non-SUSY one-loop near-miss remains.
